\section{Notation}
\label{sec:notation}

The following symbols appear throughout these notes.  Where a concept
admits more than one notation the table lists the canonical form first
and notes the section-local variant.

\begin{table}[!htbp]
\centering
\caption[General notation]{General notation used throughout these notes.}
\label{tab:notation-general}
\small
\begin{tabularx}{\textwidth}
    {l>{\raggedright\arraybackslash}p{0.34\textwidth}>{\raggedright\arraybackslash}X}
\hline
\textbf{Symbol} & \textbf{Meaning} & \textbf{Notes} \\
\hline
$z_n$        & orbit value at iteration $n$  & $z_0=0$,\; $z_{n+1}=z_n^2+c$ \\
$c$          & complex parameter (pixel)     & \\
$z_n^\star,\; c_\star$ & reference orbit / parameter
             & $z^{(0)}_n,\; c_0$ in the reuse hierarchy (\cref{subsec:ref-orbit-reuse}) \\
$\Delta z_n$ & perturbation $z_n - z_n^\star$ & \\
$\delta z_n$ & second-level perturbation      & relative to an intermediate orbit (\cref{subsec:ref-orbit-reuse}) \\
$\Delta c$   & parameter offset $c - c_\star$ & \\
$d_n \equiv \partial z_n/\partial c$ & orbit derivative
             & $d_0=0$;\; implementation starts from $d_1=1$ \\
$c^*$        & nucleus of a hyperbolic component & distinct from $c_\star$ \\
$n$          & iteration index               & \\
$p$          & period                        & \\
$n_{\text{smooth}}$ & smooth (fractional) iteration count
             & $n_{\text{smooth}} = n - \log_2\!\log_2|z_n| + \mathrm{const}$;\;
               continuous coloring \\
$N$          & generic multi-precision limb count  & section-local unless a type fixes the limb width \\
$D$          & effective number of 32-bit \code{HpSharkFloat} limbs
             & selected at runtime within a compile-time
               \code{GlobalNum\allowbreak Uint32} storage bucket \\
$L$          & number of packed base-$2^b$ coefficients
             & $L=\lceil 32D/b\rceil$ \\
$M$          & active power-of-two NTT length in the fused schedule
             & contains the complete shifted convolution support; an unshifted product needs
               $M\ge 2L-1$, rounded to a power of two, up to $2^{25}$ \\
$b$          & packed NTT coefficient width in bits
             & $b=16$ for the CUDA reference engine and its host reference model \\
\hline
\end{tabularx}
\end{table}

\begin{table}[!htbp]
\centering
\caption[Real and imaginary component notation]{Real and imaginary component notation.}
\label{tab:notation-components}
\begin{tabularx}{\textwidth}{llX}
\hline
\textbf{Symbol} & \textbf{Meaning} & \textbf{Notes} \\
\hline
$x_n,\; y_n$     & components of $z_n = x_n + i\,y_n$
                  & similarly $x_n^\star, y_n^\star$ for $z_n^\star$ \\
$\Delta x_n,\; \Delta y_n$ & components of $\Delta z_n$
                  & $\Delta c_x, \Delta c_y$ for $\Delta c$ \\
$d_{x,n},\; d_{y,n}$ & components of $d_n$
                  & \\
\hline
\end{tabularx}
\end{table}

\begin{table}[!htbp]
\centering
\caption[Linear approximation notation]{Linear approximation notation.}
\label{tab:notation-la}
\begin{tabularx}{\textwidth}{llX}
\hline
\textbf{Symbol} & \textbf{Meaning} & \textbf{Notes} \\
\hline
$z_\mathrm{ref}$ & reference orbit value at an LA entry & distinct from the global $z_n^\star$ \\
$z_j^\star$  & reference orbit value at LA stage boundary $j$
             & used when the LA hierarchy references specific stage entry points \\
$\tilde{\Delta z}$ & prepared perturbation
                  & output of the nonlinear Prepare step (\cref{sec:la-coeff-model}) \\
\code{ZCoeff},\; \code{CCoeff} & LA linear coefficients
                  & $\Delta z_\mathrm{out} \approx
                     \code{ZCoeff}\cdot\tilde{\Delta z}
                   + \code{CCoeff}\cdot\Delta c$ \\
\hline
\end{tabularx}
\end{table}

\begin{table}[!htbp]
\centering
\caption[Feature Finder notation]{Feature Finder notation.}
\label{tab:notation-feature-finder}
\begin{tabularx}{\textwidth}{llX}
\hline
\textbf{Symbol} & \textbf{Meaning} & \textbf{Notes} \\
\hline
$F(c) = z_p(c)$  & root-finding objective
                  & $p$-fold iterate of the critical point \\
$w_n$             & \code{zcoeff} product $\prod_{k=1}^{n-1} 2\,z_k$
                  & intrinsic scale of a mini-brot \\
$e_k$             & Newton iteration error at step $k$
                  & $e_k = c_k - c^*$, where $c^*$ is the target nucleus \\
\hline
\end{tabularx}
\end{table}

\begin{table}[!htbp]
\centering
\caption[Reference orbit internals notation]{Reference orbit internals notation.}
\label{tab:notation-reference-orbit}
\begin{tabularx}{\textwidth}{llX}
\hline
\textbf{Symbol} & \textbf{Meaning} & \textbf{Notes} \\
\hline
$e_n$            & reconstruction error          & $z_n^{\mathrm{HP}} - \tilde{z}_n$ \\
$\hat{z}_n$      & low-precision shadow orbit     & reduced-precision copy of $z_n^\star$ \\
$z_n^{\mathrm{HP}}$ & selected reference sample   & finite-precision source value before reduced replay \\
$z_r$            & waypoint (rebasing) reference value
                 & stored checkpoint from which intermediate iterations are replayed
                   during orbit decompression \\
$R_{\max}$       & maximum perturbation radius    & stored in \code{PerturbationResults} \\
\hline
\end{tabularx}
\end{table}

\begin{table}[!htbp]
\centering
\caption[Feature Finder internals notation]{Feature Finder internals notation.}
\label{tab:notation-feature-finder-internals}
\begin{tabularx}{\textwidth}{llX}
\hline
\textbf{Symbol} & \textbf{Meaning} & \textbf{Notes} \\
\hline
$d2_n$           & second orbit derivative        & $\partial^2 z_n / \partial c^2$;\; used in Halley iteration \\
\hline
\end{tabularx}
\end{table}
